\section{课程的历史背景}

课程随后快速回顾了语言模型的发展脉络。
从 Shannon 用语言模型估计英语熵开始，到 n-gram 在机器翻译和语音识别中的应用，再到神经语言模型、seq2seq、attention、Transformer、MoE、模型并行，最后走向 GPT、BERT、T5、GPT-3、PaLM、Chinchilla，再到近年的开源/开放权重模型。

这段历史的重点不是背名词，而是看清两个趋势：
\begin{itemize}
    \item 语言模型从“组件”变成了“核心系统”。
    \item 语言模型从“学术模型”变成了“工业基础设施”。
\end{itemize}

讲师还强调了开放性层次：
\begin{itemize}
    \item 闭源模型：只能通过 API 使用。
    \item 开放权重模型：能看权重和部分训练细节。
    \item 开源模型：权重、数据和大部分细节都开放。
\end{itemize}

这对研究很重要，因为“可复现、可检查、可分析”本身就是学习的一部分。

